\clearpage
\section{Das Vorzimmer des Universitätspräsidenten}
\label{sec:vorzimmer}
\noindent\textbf{Abschnitt:} 1.4\quad \textbf{Ort:} Verwaltungsgebäude, zweites Obergeschoss\quad
\textbf{Zeit:} 5. April 1920, vor 11 Uhr

\subsection{1.4.1 -- Ankunft im Vorzimmer}
\begin{vorlesetext}
Du folgst der Wegbeschreibung des Empfangsmitarbeiters und steigst die breite, mit dunkelrotem Teppich belegte Treppe hinauf. Das polierte Holzgeländer fühlt sich unter Deiner Hand angenehm glatt an. Mit jedem Stockwerk wird das geschäftige Treiben der Eingangshalle leiser, bis schlie\ss{}lich nur noch das gedämpfte Echo Deiner Schritte zu hören ist.

Im zweiten Obergeschoss gelangst Du in einen langen, hellen Korridor. Hohe Sprossenfenster lassen das Licht der Vormittagssonne herein, das sich in breiten Streifen auf den dunkel glänzenden Holzdielen abzeichnet. Die Wände sind bis auf Hüfthöhe mit Eichenholz verkleidet. Darüber hängen in regelmä\ss{}igen Abständen gerahmte Fotografien, historische Ansichten der Universität und einige Gemälde, die offenbar bedeutende Persönlichkeiten der Hochschule zeigen.

Mehrere schwere Holztüren führen von dem Korridor in die angrenzenden Räume. An jeder befindet sich ein kleines, sorgfältig poliertes Messingschild mit dem Namen oder der Funktion des jeweiligen Büros.

Am Ende des Ganges gelangst Du zu einer breiten Tür, deren oberes Drittel aus strukturiertem Milchglas besteht. In geschwungenen, goldfarbenen Lettern ist darauf zu lesen:

\begin{center}\textbf{Büro des Universitätspräsidenten -- Vorzimmer}\end{center}

Du betätigst den Messinggriff und trittst ein.

Der Raum ist überraschend geräumig und wirkt zugleich einladend und ausgesprochen ordentlich. Der Boden besteht aus dunklen Holzdielen, auf denen ein gro\ss{}er, fein gemusterter Teppich in gedeckten Rot- und Brauntönen liegt. Die cremefarbenen Wände sind teilweise mit Eichenholz vertäfelt. Zwei hohe Fenster bieten einen schönen Blick auf das Universitätsgelände, dessen mächtige Bäume sich sanft im Frühlingswind bewegen.

Gegenüber der Eingangstür steht ein gro\ss{}er Schreibtisch aus dunklem Eichenholz. Darauf befinden sich eine schwere Schreibmaschine, ein Telefonapparat, mehrere sorgsam gestapelte Aktenmappen sowie ein kleiner Kalender auf einem Messingständer. Hinter dem Tisch stehen hohe Aktenschränke und einige offene Regale, deren Fächer mit Briefen und Schriftstücken gefüllt sind.

An der linken Wand befindet sich eine breite, zweiflügelige Holztür. Ihre aufwendig gearbeiteten Messingbeschläge lassen vermuten, dass sich dahinter das eigentliche Arbeitszimmer des Präsidenten befindet.

Rechter Hand ist eine kleine Sitzecke eingerichtet. Vier gepolsterte Sessel gruppieren sich um einen niedrigen Tisch, auf dem eine Karaffe mit Wasser und mehrere Gläser bereitstehen. Daneben liegen einige Ausgaben der örtlichen Tageszeitung und wissenschaftliche Zeitschriften.

Drei Personen halten sich bereits in diesem Bereich auf.

Zwei Damen sitzen einander leicht zugewandt gegenüber und scheinen in ein angeregtes Gespräch vertieft zu sein. Die erste trägt zweckmä\ss{}ige, aber durchaus gepflegte Kleidung, die weniger an eine Dame der gehobenen Gesellschaft als vielmehr an jemanden erinnert, der es gewohnt ist, sich auch abseits komfortabler Salons zu bewegen. Ihre Haltung wirkt selbstbewusst und ungezwungen.

Die zweite Dame ist deutlich modischer gekleidet. Ihr Erscheinungsbild ist sorgfältig abgestimmt, und sie scheint ihrem Gegenüber mit lebhaftem Interesse zuzuhören.

Etwas abseits der beiden sitzt ein junger Mann. Er trägt einen schlichten, bereits etwas abgetragenen Anzug, dessen Stoff an den Ellbogen und Knien einen leichten Glanz angenommen hat. Seine Kleidung ist sauber und ordentlich, wirkt jedoch im Vergleich zur Garderobe der beiden Damen ausgesprochen bescheiden.

Er sitzt auffallend aufrecht, die Hände auf den Knien gefaltet, und verfolgt das Gespräch, ohne sich daran zu beteiligen. Hin und wieder wandert sein Blick zur gro\ss{}en Doppeltür auf der gegenüberliegenden Seite des Raumes. Er scheint sich in dieser Umgebung nicht sonderlich wohlzufühlen.

Hinter dem Schreibtisch erhebt sich bei Deinem Eintreten eine junge Frau, vielleicht Anfang drei\ss{}ig. Sie trägt eine schlichte, hochgeschlossene Bluse und einen dunklen Rock. Ihr Haar ist zu einer sorgfältigen Frisur hochgesteckt, und eine schmale Brille mit rundlichen Gläsern verleiht ihr ein besonders gewissenhaftes Aussehen.

Sie legt die Unterlagen, mit denen sie gerade beschäftigt war, beiseite, tritt einen kleinen Schritt hinter ihrem Schreibtisch hervor und begrü\ss{}t Dich mit einem freundlichen Lächeln.

\glqq{}Guten Morgen, mein Herr. Herzlich willkommen. Was kann ich für Sie tun?\grqq{}
\end{vorlesetext}

\subsection{1.4.2 -- Die Sekretärin des Präsidenten}
\begin{meisterinfo}
Die Sekretärin ist eine freundliche, professionelle und gut organisierte Mitarbeiterin. Sie kennt die Namen der eingeladenen Expeditionsteilnehmer und wei\ss{}, dass Professor Garmond sie um elf Uhr erwartet. Über sämtliche Einzelheiten des Auftrags ist sie nicht informiert.

\textbf{Name des NSC:} Margaret Whitmore 

\medskip\textbf{Falls Karim sich vorstellt:}

\glqq{}Herr el-Mady? Ah, selbstverständlich! Professor Garmond erwartet Sie bereits. Einen kleinen Augenblick bitte.\grqq{}

Sie gleicht seinen Namen mit einer Liste ab.

\glqq{}Wunderbar, damit wären nun alle vier Teilnehmer eingetroffen. Wenn Sie noch einen kleinen Moment Geduld hätten, werde ich dem Professor Ihre Ankunft mitteilen.\grqq{}

Sie deutet auf die Sitzecke:

\glqq{}Die anderen Herrschaften warten bereits dort. Darf ich Sie bitten, inzwischen ebenfalls Platz zu nehmen? Sie haben sicherlich Gelegenheit, sich ein wenig miteinander bekannt zu machen.\grqq{}

\medskip\textbf{Falls Karim nach den Wartenden fragt:}

\glqq{}Das sind Fräulein Collins, Fräulein Johnson und Herr Moss. Soweit ich informiert bin, wurden Sie alle zu derselben Besprechung eingeladen. Professor Garmond wird Ihnen gleich alles Weitere erläutern.\grqq{}

\medskip\textbf{Falls Karim nach dem Gegenstand der Besprechung fragt:}

\glqq{}Ich fürchte, dazu kann ich Ihnen keine näheren Auskünfte geben. Professor Garmond möchte die Angelegenheit persönlich mit Ihnen besprechen.\grqq{}

\medskip\textbf{Falls Karim fragt, ob der Präsident noch beschäftigt sei:}

\glqq{}Der Professor bereitet noch einige Unterlagen vor. Ich werde ihm gleich mitteilen, dass seine Gäste vollständig eingetroffen sind.\grqq{}

\medskip\textbf{Weiterer Ablauf:} Die Sekretärin nimmt eine Mappe, klopft zweimal an die Doppeltür, öffnet nach einem kaum verständlichen \glqq{}Herein!\grqq{} und verschwindet im Arbeitszimmer. Die Tür schlie\ss{}t sich. Nun können sich die Ermittler kennenlernen.
\end{meisterinfo}

\subsection{1.4.3 -- Das erste Zusammentreffen der Ermittler}
\begin{vorlesetextalle}
Nachdem die Sekretärin hinter der gro\ss{}en Doppeltür verschwunden ist, kehrt für einen Moment Ruhe in das Vorzimmer ein.

Durch die leicht geöffneten Fenster dringen das ferne Stimmengewirr der Studenten und das Zwitschern einiger Vögel herein. Von irgendwo im Gebäude ist das rhythmische Klappern einer Schreibmaschine zu hören.

Ansonsten herrscht eine beinahe feierliche Stille, die nur vom leisen Ticken einer Wanduhr unterbrochen wird.

Ihr befindet Euch nun zu viert in der kleinen Sitzecke. Offenbar seid Ihr alle aus unterschiedlichen Gründen zu dieser Besprechung eingeladen worden.

Noch kennt Ihr weder die genauen Absichten des Universitätspräsidenten noch die Aufgaben, die Euch erwarten.

Vorerst bleibt Euch nichts anderes übrig, als die wenigen Minuten des Wartens zu nutzen.
\end{vorlesetextalle}

\begin{meisterinfo}
\textbf{Kennenlernphase:} bewusst ohne Zeitdruck oder festgelegten Gesprächsverlauf spielen. Die Spieler entscheiden, wie ihre Figuren sprechen, reagieren und sich vorstellen.

\begin{itemize}
\item \textbf{Francis Collins (Nicole):} selbstbewusste, erfahrene Entdeckerin mit praktischem Auftreten.
\item \textbf{Lillie Johnson (Sisi):} gepflegte, kommunikative Reporterin mit Interesse an ihren Mitreisenden.
\item \textbf{Karim el-Mady (Jörg):} zuletzt eingetroffener Auslandskorrespondent.
\item \textbf{Natheniel Moss (Manuel):} einfach gekleideter Bibliothekar, der eingeschüchtert wirkt und sich zurückhält.
\end{itemize}

Wenn das Gespräch nicht von selbst in Gang kommt, sind die Berufe der Anwesenden, ihre Einladungen, frühere Expeditionen und Erwartungen an Ägypten mögliche Anknüpfungspunkte. Natheniels Zurückhaltung kann Anlass für Fragen der anderen sein.

\textbf{Dauer:} als Orientierung etwa fünf bis zehn Minuten tatsächliche Spielzeit, je nach Gesprächsverlauf.
\end{meisterinfo}

\subsection{1.4.4 -- Die Einladung zum Präsidenten}
\begin{vorlesetextalle}
Nach einigen Minuten hört Ihr hinter der schweren Doppeltür gedämpfte Stimmen. Kurz darauf erklingt das Geräusch eines sich bewegenden Stuhls, gefolgt von einigen Schritten.

Die Tür öffnet sich, und die Sekretärin tritt wieder in das Vorzimmer.

Sie schlie\ss{}t die Tür behutsam hinter sich und wendet sich mit einem freundlichen, beinahe feierlichen Lächeln an Euch.

\glqq{}Meine Damen, meine Herren, Professor Garmond ist nun bereit, Sie zu empfangen.\grqq{}

Sie tritt einen Schritt zur Seite und öffnet beide Flügel der gro\ss{}en Eichentür.

\glqq{}Wenn Sie mir bitte folgen würden. Der Professor erwartet Sie.\grqq{}
\end{vorlesetextalle}

\begin{meisterinfo}
\textbf{Szenenabschluss:} Abschnitt 1.4 endet mit der Einladung. Die Beschreibung des Arbeitszimmers sowie das erste Auftreten von Richard Garmond folgen erst in Abschnitt 1.5.

\textbf{Redaktionelle Hinweise:} Kleidung und Auftreten der Spielerfiguren sind erste Eindrücke, keine Vorgaben für ihr Rollenspiel. Die Sekretärin im Vorzimmer ist eine andere Person als der Mitarbeiter am Empfang im Erdgeschoss.
\end{meisterinfo}
